\documentclass{article}
\usepackage{notes}

\setcounter{section}{0}
\renewcommand{\thesection}{3.\arabic{section}}

\begin{document}
\stdtitle{ANALYSIS II}{Exercise Sheet 3}{Jakob Schildhauer}

\section{Lipschitz or not}
\begin{enumerate}[label=(\alph*)]
    \item $x \mapsto \frac{x + 1}{2}$
    \item $R^2 \ni (x, y) \mapsto (x + 1, y) \in \R^2$
\end{enumerate}

\section{Problems at the origin}
Let $y_k := f\left(\frac{1}{k}, \frac{1}{k}\right)$ and $z_k := f\left(\frac{1}{k}, -\frac{1}{k}\right)$ for $k \in \N$. Then $y_k \to 1$ and $z_k \to -1$. By continuity, $f(0, 0) = 1$ and $f(0, 0) = -1$. Contradiction. \\
$g$ is trivially continuous except at the origin. Let $x_k$ be any sequence with $x_k \to 0$. Then $g(x_k) \to 0$, since $|g(x_k)| \leq \abs{x_k}$. By continuity, with $g(0) = 0$, $g$ is continuous, and since the limit is unique, $g$ is uniquely defined. 


\section{Open, closed, complete and compact}
\begin{enumerate}
    \item Not open, not closed, not complete, not compact since it is partially open and partially open and you can contstruct $E_1 \ni (0, 0, k^{-1}) \to 0 \notin E_1$. 
    \item Not open, closed, complete, not compact since the preimage is closed and the set is unbounded.
    \item Open, not closed, not complete, not compact since this is the union of open sets and it is unbounded. 
    \item Open, not closed, not complete, not compact since the preimage is open. 
    \item Open, not closed, not complete, not compact since the determinant is a continuous function that maps the matrices to the open set $\R \setminus \{0\}$. 
    \item Closed?
    \item Not open, closed, complete, compact since the set is finite. 
    \item Not open, closed, complete, compact since this set is just a closed box, which was shown to be compact. 
    \item Not open, not closed, not complete since it inherits the features of both sets as the dimensions are independent. 
    \item Closed, not open, complete, compact since the intersection of compact sets is compact. 
\end{enumerate}

\newpage

\section{Multiple choice}
Correct: \titlefont{(c)}, \titlefont{(d)}, \titlefont{(e)}. 

\section{Multiple choice}
Correct: \titlefont{(a)} (if the desk is in Zürich), \titlefont{(c)}, \titlefont{(e)}. 

\section{Multiple choice}
Correct: \titlefont{(a)}, \titlefont{(c)}, \titlefont{(d)}. 

\section{Cauchy-Schwarz}
\dots

\section{Inner product and norm}
$\norm{v + w}^2 = \langle v + w, v + w\rangle 
= \langle v, v\rangle^2 + 2\langle v, w\rangle + \langle w, w\rangle^2
\leq \norm{v}^2 + 2\norm{v}\norm{w} + \norm{w}^2 = (\norm{v} + \norm{w})^2$. \\
Since the norm is non-negative, taking the square root yields the triangle inequality.

\section{All norms are equivalent in $\R^n$}
\begin{enumerate}
    \item $f(x) \leq \sum_{i=1}^n |x_i| f(e_i) \leq C_1 \norm{x}$
    \item $\abs{f(x) - f(y)} = \norm{f(x) - f(y)} \leq C \norm{x - y}$, so $f$ is Lipschitz.
    \item Since $x \neq 0$ and the sphere is compact, it attains a minimum $c_2 > 0$.  
    \item $x \neq 0 \implies f(x) = \norm{x} f\left(\frac{x}{\norm{x}}\right) \geq c_2 \norm{x}, x = 0 \implies f(0) = 0$. 
    \item Let $C := \max \Set{C_1, c_2^{-1}}$. 
\end{enumerate}

\section{Hilbert Schmidt norm of the composition}
Viewing the matrix product as a sum inner products of rows and columns, applying Cauchy-Schwarz and the definition of the Hilbert Schmidt norm, we get
\begin{align*}
    \norm{\Psi \Phi}^2 &= \sum_{i} \sum_{j} \left(\sum_{k=1}^n \Psi _{ik} \Phi_{kj}\right)^2 \\
    &\leq \sum_{i} \sum_{j} \left(\norm{\Psi _{i\cdot}}^2 \norm{\Phi_{\cdot j}}^2\right) \\
    &= \sum_{i} \norm{\Psi _{i\cdot}}^2 \sum_{j} \norm{\Phi_{\cdot j}}^2 \\
    &= \norm{\Psi }^2 \norm{\Phi}^2
\end{align*}

\section{Uniform continuity and moduli of continuity}
\begin{enumerate}
    \item Suppose $f$ has a modulus of continuity. Then, $\abs{f(x) - f(y)} \leq \omega(\abs{x - y}) \leq \omega(\delta)$ for all $x, y$ with $\abs{x - y} < \delta$. Hence, $f$ is uniformly continuous. \\
    Conversly, suppose $f$ is uniformly continuous. Then, for $\varepsilon = 1/n$, there exists $\delta_n > 0$ such that $\abs{f(x) - f(y)} < 1/n$ for all $x, y$ with $\abs{x - y} < \delta_n$. Define $\omega(\varepsilon) := 1/n$ for $\varepsilon \in [\delta_{n+1}, \delta_n)$. Then, $\omega$ is a modulus of continuity for $f$.
    \item \dots 
    \item \dots
\end{enumerate}
\end{document}